% Mục 2 (bản rút gọn). Rút từ sections/platform.tex; nguồn: results/<máy>/machine.txt, campaign.log.
Đề yêu cầu số nền tảng shared-memory không ít hơn số thành viên. Nhóm có ba thành viên và ba máy
(Bảng~\ref{tab:cross:machines}). Mục Kết quả và phân tích dựa trên \PCname{apple-m5}. RQ6 so ba máy.

Máy \PCname{apple-m5} là một Apple M5 với 10 core, gồm 4 P-core và 6 E-core, không có SMT
(simultaneous multithreading, nhiều luồng phần cứng chạy chung một core). Mỗi cụm có cache L2 riêng,
16~MiB cho cụm P-core và 6~MiB cho cụm E-core. Khi số thread vượt 4, một phần thread chạy trên E-core
chậm hơn, và loại core của mỗi thread do scheduler của macOS quyết định. AArch64 có memory model yếu hơn
x86-TSO (total store order), nên lỗi chọn sai memory order có thể quan sát được trên máy này. x86 chỉ
cho một phép đọc vượt lên trước phép ghi đứng trước nó. AArch64 cho đổi chỗ hầu hết các cặp đọc ghi tới
địa chỉ khác nhau. CPU có các lệnh atomic LSE
(Large System Extensions) của ARMv8.1, nên \texttt{fetch\_add} và \texttt{compare\_exchange} là một lệnh
máy (\texttt{ldadd}, \texttt{cas}) chứ không phải vòng LL/SC. Hệ điều hành báo cache line dài 128 byte.
Mã nguồn đệm mọi biến dùng chung theo hằng \texttt{pc::kCacheLine}, bằng 128 trên Apple arm64 và 64 trên
nền tảng khác.
% nguồn: results/apple-m5/machine.txt (l2_cache: 16 MiB per 4 P cores / 6 MiB per 6 E cores; smt: no);
% include/platform.hpp.

Phần mềm gồm macOS 27.0.1 và Apple clang 21.0.0 với cờ \texttt{-std=c++20 -O2}. Apple clang không nhận
\texttt{-fopenmp}, và macOS không có \texttt{pthread\_spinlock\_t} hay \texttt{pthread\_barrier\_t}.
Vì vậy các baseline trên máy này là \texttt{std::mutex}, \texttt{pthread\_mutex\_t},
\texttt{os\_unfair\_lock}, \texttt{std::barrier} và barrier dựng từ mutex và condition variable. Đợt đo
chính dài 1 giờ 26 phút, gồm 12\,152 lần chạy hoàn thành (cùng 9 lần vượt thời hạn). Máy cắm sạc và chế độ
tiết kiệm điện tắt. Runner chờ tới khi load average một phút xuống dưới 7,0 mới bắt đầu (6,73). Load average là số thread
sẵn sàng chạy, trung bình trong một phút, do hệ điều hành báo.
% nguồn: campaign.log (START load1=6.73; END duration=1:25:45 ok=12152 timeout=9); machine.txt: power AC
% Power, low_power_mode 0.

Máy \PCname{C4M7Y670NT} là một Apple M3 Pro với 12 core, gồm 6 P-core và 6 E-core, không có SMT, với
cache L2 16~MiB cho cụm P-core và 4~MiB cho cụm E-core. Máy chạy macOS 26.7.1 và Apple clang 21.0.0 với
cùng cờ biên dịch và cùng các baseline như \PCname{apple-m5}. Lưới thread thêm 24. Đợt đo dài 1 giờ 53
phút, gồm 13\,244 lần chạy hoàn thành (cùng 10 lần vượt thời hạn), máy cắm sạc.
% nguồn: results/C4M7Y670NT/machine.txt; campaign.log (END duration=1:53:10 ok=13244 timeout=10).

Máy \PCname{LAP1184} là một laptop Intel Core i5-1335U. Theo thông số của Intel, CPU này có 2 P-core với
Hyper-Threading và 8 E-core, tức 10 core và 12 luồng phần cứng. Nhóm đo trong Ubuntu 20.04 chạy trên
WSL2, một máy ảo Hyper-V của Windows 11. Máy ảo có 12 CPU ảo và báo chúng là 6 core, mỗi core 2 luồng,
không khớp với CPU thật. x86-64 có memory model TSO và cache line 64 byte. Phần mềm gồm g++ 11.5 với
\texttt{-std=c++20 -O2 -fopenmp}. Ngoài các baseline chung, máy này có \texttt{pthread\_spinlock\_t},
\texttt{pthread\_barrier\_t}, \texttt{omp\_lock\_t} và \texttt{\#pragma omp barrier}, nên có 13 lock, 9
barrier và 20 counter. Đợt đo dài 1 giờ 20 phút 36 giây, gồm 14\,322 lần chạy hoàn thành (cùng 12 lần
vượt thời hạn).
% nguồn: results/LAP1184/machine.txt (virtualization: Microsoft; physical_cores: 6; threads_per_core: 2);
% campaign.log (END duration=1:20:36 ok=14322 timeout=12).
